\chapter{Linear spanners}
\label{chap:linear-spanners}

\noindent Formalized in \texttt{TemporalGraph/LinearSpanner/ExtremallyMatched.lean}
(extremally matched bi-cliques), \texttt{Reduction.lean} (dismountability in
bi-cliques, Lemma~5), \texttt{Stars.lean} (the edge sets of the stars),
\texttt{SimpleStar.lean} (Observation~9), \texttt{ExtendedStar.lean}
(Lemmas~13 to~16), \texttt{Bound.lean} (Lemma~17, Theorem~18 and the
bi-clique Theorem~3) and \texttt{Clique.lean} (Theorems~20 and~2), with the
supporting API in \texttt{Basic.lean}, \texttt{Connectivity/Basic.lean},
\texttt{TemporalWalk.lean} and \texttt{Dismountability/OneSided*.lean}.
This chapter follows Júlia Baligács, \emph{Temporal Cliques Admit Linear
Spanners}, arXiv:2606.05156v1 (2026) --- abbreviated Baligács below, an arXiv
preprint only --- which answers the question of Casteigts, Peters \&
Schoeters: temporal cliques admit spanners of size $O(n)$, improving the
$O(n \log n)$ bound of Chapter~\ref{chap:main-theorem}. The numbering below is
that of the arXiv version v1, which uses a single counter for all statements
(Theorem~2, Lemma~5, Definition~8, Observation~9, \ldots).

Design choices that differ from the paper.
\begin{itemize}
\item
  The paper assumes without loss of generality (\S~2.1) that every edge
  carries a single label and that no two edges at a vertex share a label.
  Here these are hypotheses: the graphs are simple
  (\texttt{IsSimple}, Definition~\ref{def:is-simple-clique}) and, where the
  paper uses the uniqueness of extremal edges, proper
  (\texttt{IsProper}, Definition~\ref{def:is-proper-happy}).
\item
  No index is ever built. The paper compares positions
  $\mathrm{pos}_v(u) < \mathrm{pos}_v(w)$ in the order a vertex $v$ induces on
  its neighbors, and numbers the vertices along the order induced by a center
  (the $s^*$- and $t^*$-ordered labelings, Definition~7). Here
  $\mathrm{pos}_v(u) \le \mathrm{pos}_v(w)$ is the comparison of arc times
  $\mathrm{arcTime}(v, u) \le \mathrm{arcTime}(v, w)$, and ``$j \le i$'' in a
  labeling centered at $s^*$ is a comparison of the times of the edges at
  $s^*$. A cut at index $k$ (Definition~8) is given by a threshold time $\tau$.
\item
  Absence is explicit. Arc times are \texttt{arcTime : Option ℕ}, and $N_{\min}$,
  $N_{\max}$ are the functions \texttt{earliestNeighborFn},
  \texttt{latestNeighborFn} with values in \texttt{Option V}, \texttt{none}
  standing for the absence of an edge (no junk value). A comparison is thus
  stated as ``some value of $\mathrm{arcTime}(u, v)$ is at most some value of
  $\mathrm{arcTime}(u, w)$'', false when one of the arcs is absent.
\item
  A spanner of a bi-clique is a subgraph $\mathrm{fromEdgeSet}(E)$ of an edge
  set $E$ (Definition~\ref{def:from-arc-edge-set}), and a bi-clique spanner
  is a one-sided spanner (Definition~\ref{def:spanner-from-to}): there is no
  separate ``covering'' predicate.
\end{itemize}
\bigskip

\section*{Preliminaries (Baligács \S~2.1)}

\begin{definition}[Arc times --- Baligács \S\S~1, 2.1]
  \label{def:arc-time}
  \uses{def:is-simple-clique,def:temporal-graph}
  \lean{TemporalGraph.arcTime}
  \lean{TemporalGraph.arcTime_eq_some}
  \lean{TemporalGraph.arcTime_comm}
  \leanok
  Let $G$ be a temporal graph. $\mathrm{arcTime}(u, v)$ is the time of an arc
  from $u$ to $v$, the paper's $\lambda(\{u, v\})$, as an
  \texttt{Option ℕ}: \texttt{none} when $u$ and $v$ are not adjacent, and
  \texttt{some n} for the time $n$ of the arc when $G$ is simple. When $G$ is
  moreover symmetric, $\mathrm{arcTime}(u, v) = \mathrm{arcTime}(v, u)$.
\end{definition}

\begin{definition}[Reaching within a time window --- Baligács \S~2.1]
  \label{def:reachable-within}
  \uses{def:reachable,def:temporal-walk}
  \lean{TemporalGraph.ReachableWithin}
  \lean{TemporalGraph.ReachableBefore}
  \lean{TemporalGraph.ReachableWithin.step}
  \leanok
  $x$ \emph{reaches $y$ within $[a, b]$} when some temporal walk, possibly
  empty, goes from $x$ to $y$, leaves $x$ at a time at least $a$ and arrives at
  $y$ at a time at most $b$; $x$ \emph{reaches $y$ before $b$} is the window
  $[0, b]$. A walk is extended by one arc at its end
  (\texttt{ReachableWithin.step}): the paper's temporal walks, whose repeated
  vertices are collapsed, are built with it one edge at a time.
\end{definition}

\begin{lemma}[Edges of a walk --- Baligács \S~2.1]
  \label{lem:walk-edges}
  \uses{def:reachable-within,def:from-arc-edge-set}
  \lean{TemporalGraph.TemporalWalk.edges}
  \lean{TemporalGraph.TemporalWalk.ncard_setOf_mem_edges_le}
  \lean{TemporalGraph.TemporalWalk.reachableBefore_fromEdgeSet}
  \leanok
  A temporal walk $p$ traverses the list $\mathrm{edges}(p)$ of edges (an
  unoriented edge together with its time), of which there are at most
  $\mathrm{length}(p)$ distinct ones. If all edges of $p$ belong to a set $E$,
  the walk lives in $\mathrm{fromEdgeSet}(E)$: the endpoints of $p$ are
  reachable one from the other there, before the arrival time of $p$.
\end{lemma}

\begin{proof}
  \leanok
  By induction on the walk.
\end{proof}

\section*{Dismountability and extremally matched bi-cliques (Baligács \S~2.3)}

\begin{definition}[Dismountable sources and targets --- Baligács \S~2.3]
  \label{def:source-target-dismountable}
  \uses{def:is-biclique,def:earliest-latest-neighbor,def:spanner-from-to}
  \lean{TemporalGraph.SourceDismountable}
  \lean{TemporalGraph.OneSidedDismountable}
  \leanok
  Let $G$ be a temporal graph with sources $S$ and targets $T$. A source
  $x \in S$ is \emph{dismountable} when another source $x'$ has its earliest
  edge to a target $y$, at time $n$, and the edge $yx$ occurs at a time
  $m \le n$: $x$ reaches through $x \to y \to x'$ everything $x'$ reaches.
  Dually, a target $x \in T$ is dismountable when another target $x'$ has its
  latest edge to a source $y$, at time $n$, and the edge $yx$ occurs at a time
  $m \ge n$. Target dismountability is the one-sided dismountability of
  Definition~\ref{def:spanner-from-to}
  (\texttt{OneSidedDismountable}), and source dismountability
  (\texttt{SourceDismountable}) is its mirror image.
\end{definition}

\begin{theorem}[Mounting back a source or a target --- Baligács \S~2.3]
  \label{thm:source-dismountability}
  \uses{def:source-target-dismountable,def:spanner-from-to,thm:one-sided-dismountability}
  \lean{TemporalGraph.Subgraph.sourceDismountability}
  \lean{TemporalGraph.Subgraph.isSpannerFromTo_of_sourceDismount}
  \leanok
  Let $G$ be symmetric and $x$ a dismountable source of a bi-clique subgraph
  $H$ on $S$ and $T$. Every symmetric one-sided spanner of $H$ without $x$
  extends to a symmetric one-sided spanner of $H$ from $S$ to $T$ with two more
  edges $x'y$ and $yx$. Dually for a dismountable target
  (Theorem~\ref{thm:one-sided-dismountability}).
\end{theorem}

\begin{proof}
  \leanok
  \uses{lem:reachable-of-earliestEdge,lem:reachable-of-latestEdge}
  The earliest edge $x'y$ of $H$ stays the earliest edge at $x'$ in the smaller
  graph. So $x \to y \to x'$, at times $m \le n$, followed by any walk from $x'$
  is a temporal walk (Lemma~\ref{lem:reachable-of-earliestEdge}), and $x$
  reaches every target that $x'$ reaches. The target case appends
  $x' \to y \to x$ to any walk into $x'$ (Lemma~\ref{lem:reachable-of-latestEdge}).
\end{proof}

\begin{definition}[Extremally matched bi-clique --- Baligács \S~2.3]
  \label{def:extremally-matched}
  \uses{def:is-biclique,def:earliest-latest-neighbor}
  \lean{TemporalGraph.IsExtremallyMatched}
  \leanok
  A temporal graph $G$ is an \emph{extremally matched} bi-clique on sources $S$
  and targets $T$ if it is a bi-clique on $S$ and $T$ in which the earliest
  edges form a perfect matching $E_{\min}$ and the latest edges a perfect
  matching $E_{\max}$: the earliest edge at any vertex of $S \cup T$ is also the
  earliest edge at its other endpoint, and likewise for the latest edges.
\end{definition}

\begin{lemma}[Basic properties of extremally matched bi-cliques --- Baligács \S~2.3]
  \label{lem:extremally-matched-basic}
  \uses{def:extremally-matched,def:neighbor-functions,lem:neighbor-fn-eq,def:is-proper-happy}
  \lean{TemporalGraph.IsExtremallyMatched.earliestNeighborFn_eq_some_symm}
  \lean{TemporalGraph.IsExtremallyMatched.ncard_eq}
  \leanok
  In a proper extremally matched bi-clique, if $N_{\min}(v) = \mathrm{some}\ w$
  then $N_{\min}(w) = \mathrm{some}\ v$ (and likewise for $N_{\max}$), and
  $|S| = |T|$.
\end{lemma}

\begin{proof}
  \leanok
  The earliest edge $vw$ at $v$ is also the earliest edge at $w$, so by
  uniqueness in a proper graph $N_{\min}(v) = \mathrm{some}\ w$ and
  $N_{\min}(w) = \mathrm{some}\ v$; the same holds for $N_{\max}$. Hence
  $N_{\min}$ is injective on $S \cup T$, and it maps $S$ onto $T$, so
  $|S| = |T|$.
\end{proof}

\begin{theorem}[Temporal connectivity of a spanner of an extremally matched bi-clique --- Baligács \S~2.3]
  \label{thm:extremally-matched-tc}
  \uses{def:extremally-matched,def:spanner-from-to,thm:tc-of-spanner-from-to}
  \lean{TemporalGraph.IsExtremallyMatched.temporallyConnected}
  \leanok
  Let $H$ be an extremally matched bi-clique on $S$ and $T$, and $K$ a
  symmetric subgraph containing the matchings $E_{\min}$ and $E_{\max}$ and in
  which every source reaches every target. Then $K$ is temporally connected on
  $S \cup T$.
\end{theorem}

\begin{proof}
  \leanok
  \uses{thm:tc-of-spanner-from-to}
  Since $K$ contains $E_{\min}$ and $E_{\max}$, a target reaches whatever its
  $N_{\min}$ reaches and a source is reached by whatever reaches its $N_{\max}$.
  This is Theorem~\ref{thm:tc-of-spanner-from-to} (SAND 2025, \S~4.1).
\end{proof}

\begin{lemma}[No dismountable vertex --- Baligács \S~2.3]
  \label{lem:extremally-matched-of-not-dismountable}
  \uses{def:source-target-dismountable,def:extremally-matched}
  \lean{TemporalGraph.Subgraph.isExtremallyMatched_of_not_dismountable}
  \leanok
  A finite, simple, symmetric bi-clique subgraph with two nonempty parts and
  without dismountable sources or targets is extremally matched.
\end{lemma}

\begin{proof}
  \leanok
  Let $v$ be a source, $w = N_{\min}(v)$ and $u = N_{\min}(w)$. If $u \ne v$,
  then $u$ is dismountable via $(v, w)$, since its edge to $w$ is no later than
  the earliest edge of $v$; so $N_{\min}(N_{\min}(v)) = v$ and $|S| \le |T|$.
  Dually, $N_{\max}(N_{\max}(v)) = v$ for every target $v$, and $|T| \le |S|$.
  Both maps are then bijections between $S$ and $T$, which gives the two
  perfect matchings.
\end{proof}

\begin{lemma}[Dismountability --- Baligács Lemma 5]
  \label{lem:bal-dismountability}
  \uses{def:extremally-matched,def:between,def:spanner-from-to,thm:source-dismountability,lem:extremally-matched-of-not-dismountable,lem:extremally-matched-basic}
  \lean{TemporalGraph.Subgraph.exists_isExtremallyMatched_reduction}
  \leanok
  Let $H$ be a finite, simple, symmetric bi-clique subgraph on $S$ and $T$.
  There are $S' \subseteq S$ and $T' \subseteq T$ such that the bi-clique
  $H[S' \cup T']$ induced on them is extremally matched on $S'$ and $T'$, and
  every symmetric one-sided spanner $K$ of $H[S' \cup T']$ from $S'$ to $T'$
  extends to a symmetric one-sided spanner $K'$ of $H$ from $S$ to $T$ with
  \[
    |K'| + 2(|S'| + |T'|) \le |K| + 2(|S| + |T|),
  \]
  that is, with at most $2(|S| + |T| - |S'| - |T'|)$ more edges. The bound is
  additive, which avoids truncated subtraction, and no properness is needed.
\end{lemma}

\begin{proof}
  \leanok
  By strong induction on $|S| + |T|$. If a part is empty, take
  $S' = T' = \emptyset$ and $K' = K$. If some vertex $x$ is dismountable,
  recurse on the bi-clique without $x$ and add the two edges of
  Theorem~\ref{thm:source-dismountability}, two edges for one vertex less.
  Otherwise $H$ itself is extremally matched
  (Lemma~\ref{lem:extremally-matched-of-not-dismountable}): take $S' = S$,
  $T' = T$ and $K' = K$.
\end{proof}

\section*{Cuts in stars (Baligács \S~3)}

\begin{definition}[Simple star and cut-crossing sources --- Baligács Definition 8]
  \label{def:simple-star-cut-crossing}
  \uses{def:extremally-matched,def:arc-time,def:neighbor-functions,def:temporal-walk,def:incidence-set}
  \lean{TemporalGraph.eMin}
  \lean{TemporalGraph.eMax}
  \lean{TemporalGraph.incidentEdges}
  \lean{TemporalGraph.simpleStar}
  \lean{TemporalGraph.IsCutCrossing}
  \leanok
  Let $G$ be an extremally matched bi-clique and $s^* \in S$. The matchings
  $E_{\min}$ and $E_{\max}$ are the sets \texttt{eMin}$(S)$ and
  \texttt{eMax}$(S)$ of the edges $\{v, N_{\min}(v)\}$ and $\{v, N_{\max}(v)\}$
  for $v \in S$, at their times.
  \begin{enumerate}
  \item[(i)] The \emph{simple star} $\mathrm{Star}(s^*)$ consists of $E_{\min}$
    and of the edges at $s^*$ (\texttt{incidentEdges}).
  \item[(ii)] In the $s^*$-ordered labeling, where the targets
    $t_0, \ldots, t_{n-1}$ are sorted by the times of their edges to $s^*$ and
    $s_i = N_{\min}(t_i)$, the index of a target $u$ is read on
    $\mathrm{arcTime}(s^*, u)$ and that of a source $v$ on
    $\mathrm{arcTime}(s^*, N_{\min}(v))$. For a cut given by a threshold
    $\tau$, a source $v$ is \emph{$m$-hop-cut-crossing} if its index is above
    the cut ($\tau$ is less than the time of $s^*N_{\min}(v)$) and some
    temporal walk of length at most $m$ goes from $v$ to a target
    $t_j = N_{\min}(s_j)$ whose edge to $s^*$ has time at most $\tau$, ending
    no later than the edge $\{t_j, s^*\}$.
  \end{enumerate}
\end{definition}

\begin{lemma}[Simple stars cover the cut --- Baligács Observation 9]
  \label{lem:reachable-simple-star}
  \uses{def:simple-star-cut-crossing,lem:walk-edges}
  \lean{TemporalGraph.IsExtremallyMatched.exists_reachable_simpleStar}
  \leanok
  Let $H$ be a simple proper extremally matched bi-clique subgraph, $s^* \in S$,
  $\tau$ a threshold and $m$ a length, and let $S_{\mathrm{cross}}$ be the
  $m$-hop-cut-crossing sources. There is a set $E$ of at most
  $m\,|S_{\mathrm{cross}}|$ edges of $H$ such that, in
  $\mathrm{fromEdgeSet}(E \cup \mathrm{Star}(s^*))$, every source that is
  early (the time of $s^*N_{\min}(v)$ is at most $\tau$) or crossing reaches
  every target $u$ with $\tau \le \mathrm{arcTime}(s^*, u)$.
\end{lemma}

\begin{proof}
  \leanok
  An early source $s_j$ reaches an upper target $t$ through
  $s_j \to t_j \to s^* \to t$: the edge $s_j t_j$ is the earliest at $t_j$, and
  the times at $s^*$ increase along the labeling. A crossing source follows its
  own walk, whose at most $m$ edges form $E$
  (Lemma~\ref{lem:walk-edges}), to some lower $t_j$, then continues through
  $t_j \to s^* \to t$ using only edges of $\mathrm{Star}(s^*)$.
\end{proof}

\section*{Extended stars (Baligács \S~4)}

\begin{definition}[Extended stars --- Baligács Definition 12]
  \label{def:extended-star}
  \uses{def:simple-star-cut-crossing}
  \lean{TemporalGraph.IsSourceExtendTarget}
  \lean{TemporalGraph.IsTargetExtendSource}
  \lean{TemporalGraph.extStarSource}
  \lean{TemporalGraph.extStarTarget}
  \leanok
  Let $s^* \in S$ and $t^* \in T$.
  \begin{enumerate}
  \item[(a)] The \emph{extend target} of a vertex $v$ is the last target $u$,
    in the $s^*$-ordered labeling, whose edge to $v$ comes after its edge to
    $s^*$ (\texttt{IsSourceExtendTarget}). The \emph{extended star}
    $\mathrm{Ext}(s^*)$ (\texttt{extStarSource}) consists of $E_{\min}$,
    $E_{\max}$, all the edges at $s^*$, and the edge from each source to its
    extend target.
  \item[(b)] The \emph{extend source} of a vertex $v$ is the last source $u$, in
    the $t^*$-ordered labeling (sources by decreasing time at $t^*$, and
    $t_i = N_{\max}(s_i)$), whose edge to $v$ comes before its edge to $t^*$
    (\texttt{IsTargetExtendSource}). The extended star $\mathrm{Ext}(t^*)$
    (\texttt{extStarTarget}) consists of $E_{\min}$, $E_{\max}$, all the edges at
    $t^*$, and the edge from each target to its extend source.
  \end{enumerate}
  ``Last'' is read on arc times at the center: $u'$ is no later than $u$ when
  some value of $\mathrm{arcTime}(s^*, u')$ is at most some value of
  $\mathrm{arcTime}(s^*, u)$.
\end{definition}

\begin{lemma}[Pairs reached by extended stars --- Baligács Lemma 13]
  \label{lem:ext-star-reachable}
  \uses{def:extended-star,lem:extremally-matched-basic}
  \lean{TemporalGraph.IsExtremallyMatched.reachable_extStarSource}
  \lean{TemporalGraph.IsExtremallyMatched.reachable_extStarTarget}
  \leanok
  \begin{enumerate}
  \item[(a)] In the $s^*$-ordered labeling, let $u'$ be a target,
    $z = N_{\max}(u')$ and $u$ the extend target of $z$. For a source $v$ with
    $x = N_{\min}(v)$ whose edge to $s^*$ is no later than $u s^*$, the star
    $\mathrm{Ext}(s^*)$ reaches $u'$ from $v$.
  \item[(b)] Dually, in the $t^*$-ordered labeling, for a source $v$ and
    $\mathrm{Ext}(t^*)$ reaching the targets whose $N_{\max}$ lies no later
    than the extend source of $N_{\min}(v)$.
  \end{enumerate}
\end{lemma}

\begin{proof}
  \leanok
  For (a), the walk $(v, x, s^*, u, z, u')$ lies in $\mathrm{Ext}(s^*)$ and is
  temporal: $vx$ is the earliest edge at $x$, the times at $s^*$ increase with
  the index, the definition of the extend target orders the edges at $u$, and
  $zu'$ is the latest edge at $z$. Part (b) uses the dual walk through $t^*$.
\end{proof}

\begin{lemma}[Pairs not reached by extended stars --- Baligács Lemma 14]
  \label{lem:ext-star-not-reachable}
  \uses{lem:ext-star-reachable}
  \lean{TemporalGraph.IsExtremallyMatched.arcTime_le_of_not_reachable_extStarSource}
  \lean{TemporalGraph.IsExtremallyMatched.arcTime_le_of_not_reachable_extStarTarget}
  \leanok
  Let $v \in S$, $u' \in T$, $x = N_{\min}(v)$ and $z = N_{\max}(u')$.
  \begin{enumerate}
  \item[(a)] If $\mathrm{Ext}(s^*)$ does not reach $u'$ from $v$, then
    $\{x, z\}$ occurs no later than $\{x, s^*\}$.
  \item[(b)] If $\mathrm{Ext}(t^*)$ does not reach $u'$ from $v$, then
    $\{z, t^*\}$ occurs no later than $\{z, x\}$.
  \end{enumerate}
\end{lemma}

\begin{proof}
  \leanok
  For (a), suppose the converse inequality. Then the set of targets $y$ whose
  edge to $s^*$ is before its edge to $z$ is nonempty, as it contains $x$, so
  $z$ has an extend target, and that one comes no earlier than $x$ in the
  labeling. Lemma~\ref{lem:ext-star-reachable}~(a) then gives a walk from $v$
  to $u'$, a contradiction. Part (b) is dual.
\end{proof}

\begin{corollary}[Ordering of unreached pairs --- Baligács Corollary 15]
  \label{cor:ext-star-chain}
  \uses{lem:ext-star-not-reachable}
  \lean{TemporalGraph.IsExtremallyMatched.arcTime_chain_of_not_reachable}
  \leanok
  If $\mathrm{Ext}(s^*) \cup \mathrm{Ext}(t^*)$ does not reach $u'$ from $v$,
  then the walk $(N_{\min}(t^*), t^*, N_{\max}(u'), N_{\min}(v), s^*)$ is
  temporal. Its first junction is temporal for free, $N_{\min}(t^*)$ being the
  earliest neighbor of $t^*$; the statement gives the other two inequalities,
  at $N_{\max}(u')$ and at $N_{\min}(v)$.
\end{corollary}

\begin{proof}
  \leanok
  Neither extended star reaches $u'$ from $v$: apply both parts of
  Lemma~\ref{lem:ext-star-not-reachable}, and use the symmetry of arc times
  to align the two inequalities.
\end{proof}

\begin{lemma}[Crossing or reaching --- Baligács Lemma 16]
  \label{lem:cut-crossing-or-reachable}
  \uses{cor:ext-star-chain,def:simple-star-cut-crossing,def:extended-star}
  \lean{TemporalGraph.IsExtremallyMatched.isCutCrossing_or_reachable}
  \leanok
  In the $s^*$-ordered labeling cut at a threshold $\tau$, every source $v$
  above the cut either is $3$-hop-cut-crossing, or
  $\mathrm{Ext}(s^*) \cup \mathrm{Ext}(N_{\min}(v))$ reaches every target from
  every source $s_j$ below the cut.
\end{lemma}

\begin{proof}
  \leanok
  Suppose some pair $(s_j, u')$ with $s_j$ below the cut is not reached. By
  Corollary~\ref{cor:ext-star-chain} with $t^* = N_{\min}(v)$, whose earliest
  neighbor is $v$, the walk $(v, N_{\min}(v), N_{\max}(u'), t_j, s^*)$ is
  temporal. Its first three edges reach the lower target $t_j$ no later than
  the edge $\{t_j, s^*\}$, so $v$ is $3$-hop-cut-crossing.
\end{proof}

\begin{lemma}[Size of two extended stars --- Baligács, proof of Lemma 17]
  \label{lem:ext-star-ncard}
  \uses{def:extended-star}
  \lean{TemporalGraph.IsExtremallyMatched.ncard_extStar_union_le}
  \leanok
  $|\mathrm{Ext}(s^*) \cup \mathrm{Ext}(t^*)| \le 6n$, where $n = |S|$.
\end{lemma}

\begin{proof}
  \leanok
  The union consists of six groups of at most $n$ edges each: $E_{\min}$,
  $E_{\max}$, the edges at $s^*$, the edges at $t^*$, at most one extend edge
  per source and at most one per target (extend targets and sources are
  unique, by properness).
\end{proof}

\begin{lemma}[Median threshold]
  \label{lem:exists-threshold}
  \lean{Set.exists_median_threshold}
  \leanok
  For a finite set $A$ and any $f : V \to \mathbb{N}$, some threshold $\tau$
  leaves at most $\lfloor |A| / 2 \rfloor$ elements of $A$ with $f < \tau$ and
  at most $\lfloor |A| / 2 \rfloor$ with $f > \tau$.
\end{lemma}

\begin{proof}
  \leanok
  Take for $\tau$ the least value leaving more than $\lfloor |A| / 2 \rfloor$
  elements with $f \le \tau$; it exists since $f$ is bounded on $A$. By
  minimality, fewer elements have $f < \tau$, and the elements with $f > \tau$
  form the complement of a set of more than $\lfloor |A| / 2 \rfloor$ elements.
\end{proof}

\begin{lemma}[Reaching half of the bi-clique --- Baligács Lemma 17]
  \label{lem:reachable-half}
  \uses{lem:reachable-simple-star,lem:cut-crossing-or-reachable,lem:ext-star-ncard}
  \lean{TemporalGraph.IsExtremallyMatched.exists_reachable_of_threshold}
  \leanok
  Let $H$ be a finite, simple, proper, symmetric extremally matched bi-clique
  subgraph, $s^* \in S$, and $\tau$ an arbitrary threshold time cutting the
  $s^*$-ordered labeling. Some set $E$ of at most $6|S|$ edges of $H$,
  containing $E_{\min} \cup E_{\max}$, is such that in
  $\mathrm{fromEdgeSet}(E)$ either every source reaches every target above the
  cut, or every source below the cut reaches every target.
\end{lemma}

\begin{proof}
  \leanok
  If every source above the cut is $3$-hop-cut-crossing, Observation~9
  (Lemma~\ref{lem:reachable-simple-star}) gives $E'$ with $|E'| \le 3|S|$ such
  that $E' \cup \mathrm{Star}(s^*) \cup E_{\max}$, of size at most
  $3|S| + 2|S| + |S|$, reaches every target above the cut from every source.
  Otherwise some source $v$ above the cut is not crossing, and
  $\mathrm{Ext}(s^*) \cup \mathrm{Ext}(N_{\min}(v))$, of size at most $6|S|$
  (Lemma~\ref{lem:ext-star-ncard}), reaches every target from every source
  below the cut by Lemma~\ref{lem:cut-crossing-or-reachable}.
\end{proof}

\section*{Linear spanners of bi-cliques}

\begin{theorem}[Extremally matched bi-cliques --- Baligács Theorem 18 (and Theorem 6)]
  \label{thm:spanner-extremally-matched}
  \uses{lem:reachable-half,lem:exists-threshold,lem:bal-dismountability,lem:extremally-matched-basic,thm:extremally-matched-tc,lem:spanner-from-to-split}
  \lean{TemporalGraph.IsExtremallyMatched.exists_isSpannerFromTo_temporallyConnected}
  \leanok
  Let $H$ be a finite, simple, proper, symmetric extremally matched bi-clique
  subgraph on $S$ and $T$, with $n = |S| = |T|$. There is a symmetric
  one-sided spanner of $H$ from $S$ to $T$ with at most $14n$ edges, which is
  moreover temporally connected on $S \cup T$.
\end{theorem}

\begin{proof}
  \leanok
  By strong induction on $n$, with a spanner containing $E_{\min} \cup
  E_{\max}$. Cut the $s^*$-ordered labeling at a median threshold
  (Lemma~\ref{lem:exists-threshold}) and apply
  Lemma~\ref{lem:reachable-half}: at most $6n$ edges reach every target above
  the cut from every source, or every target from every source below the cut.
  What is left, a bi-clique with one side of size at most $\lfloor n/2 \rfloor$,
  is reduced by Lemma~\ref{lem:bal-dismountability} to an extremally matched
  bi-clique of size $x \le \lfloor n/2 \rfloor$; the induction hypothesis
  spans it with $14x$ edges, which Lemma~\ref{lem:bal-dismountability} extends
  to the whole leftover with at most $10 \lfloor n/2 \rfloor + 2(n + \lfloor n/2 \rfloor)$
  edges, for a total of at most $14n$ (the two halves are glued by
  Lemma~\ref{lem:spanner-from-to-split}). Since the spanner contains $E_{\min}$
  and $E_{\max}$, Theorem~\ref{thm:extremally-matched-tc} gives temporal
  connectivity on $S \cup T$.
\end{proof}

\begin{theorem}[Linear spanners of temporal bi-cliques --- Baligács Theorem 3]
  \label{thm:linear-biclique-spanner}
  \uses{lem:bal-dismountability,thm:spanner-extremally-matched,def:spanner-from-to}
  \lean{TemporalGraph.Subgraph.exists_isSpannerFromTo_ncard_le_linear}
  \leanok
  Let $G$ be symmetric and loopless and $H$ a finite, simple, proper,
  symmetric temporal bi-clique subgraph on $S$ and $T$. Then $H$ has a
  symmetric one-sided spanner from $S$ to $T$
  (Definition~\ref{def:spanner-from-to}) with at most
  $10 \min(|S|, |T|) + 2(|S| + |T|)$ edges, to be compared with the
  $O(n \log n)$ bound of Theorem~\ref{thm:exists-spanner-from-to-ncard-le}.
\end{theorem}

\begin{proof}
  \leanok
  \uses{thm:exists-spanner-from-to-ncard-le}
  Lemma~\ref{lem:bal-dismountability} gives $S'$, $T'$ with
  $x = |S'| = |T'| \le \min(|S|, |T|)$ (Lemma~\ref{lem:extremally-matched-basic})
  and an extremally matched bi-clique, which Theorem~\ref{thm:spanner-extremally-matched}
  spans with at most $14x$ edges. The extension of Lemma~\ref{lem:bal-dismountability}
  has $|K'| + 2(|S'| + |T'|) \le 14x + 2(|S| + |T|)$, that is
  $|K'| \le 10x + 2(|S| + |T|)$.
\end{proof}

\section*{Linear spanners of temporal cliques (Baligács \S~5)}

\begin{theorem}[Non-dismountable cliques give extremally matched bi-cliques --- Baligács Theorem 20]
  \label{thm:extremally-matched-of-not-two-hop}
  \uses{def:extremally-matched,def:two-hop-dismountable,def:vminus-vplus-vzero,def:between,cor:biclique-of-not-two-hop,lem:reciprocity,lem:vminus-vplus-matchings,lem:earliestEdge-exists-finite-clique}
  \lean{TemporalGraph.Subgraph.isExtremallyMatched_between_of_not_twoHopDismountable}
  \leanok
  Let $H$ be a finite, nontrivial, simple, proper, symmetric clique that is not
  $\{1,2\}$-hop dismountable. Then the restriction
  $H.\mathrm{between}(V^-, V^+)$ of $H$ to the edges between $V^-$ and $V^+$ is
  an extremally matched bi-clique with sources $V^+$ and targets $V^-$. This is
  the paper's Theorem~20, quoted there from SAND 2025, Corollary~4.1.
\end{theorem}

\begin{proof}
  \leanok
  The restriction is a bi-clique by Corollary~\ref{cor:biclique-of-not-two-hop}.
  The earliest edge $va$ of $H$ at $v \in V^+$ goes to $a \in V^-$
  (Lemma~\ref{lem:vminus-vplus-matchings}), and by reciprocity
  (Lemma~\ref{lem:reciprocity}, SAND 2025, Lemma~3.5) the edges from $a$ to
  $V^+$ occur no earlier: it is the earliest edge at both ends in the
  bi-clique. Every vertex of $V^-$ ends such an edge, which gives $E_{\min}$;
  $E_{\max}$ is dual.
\end{proof}

\begin{theorem}[Spanner of a non-dismountable clique, linear size --- Baligács, proof of Theorem 2]
  \label{thm:spanner-of-not-two-hop-linear}
  \uses{thm:extremally-matched-of-not-two-hop,thm:spanner-extremally-matched,lem:vzero-empty,lem:card-bounds,def:spanner}
  \lean{TemporalGraph.Subgraph.exists_isSpannerOf_ncard_le_seven_mul_of_not_twoHopDismountable}
  \leanok
  Let $G$ be symmetric and loopless. A finite, nontrivial, simple, proper,
  symmetric clique $H$ that is not $\{1,2\}$-hop dismountable has a symmetric
  spanner with at most $14|V^+| = 7|\mathrm{Verts}(H)|$ edges.
\end{theorem}

\begin{proof}
  \leanok
  Theorem~\ref{thm:spanner-extremally-matched} applied to the extremally matched
  bi-clique of Theorem~\ref{thm:extremally-matched-of-not-two-hop} gives a
  one-sided spanner with at most $14|V^+|$ edges, temporally connected on
  $V^+ \cup V^-$. Since $V^0 = \emptyset$ (Lemma~\ref{lem:vzero-empty}), this
  is all of $\mathrm{Verts}(H)$, so it is a spanner of the clique $H$, and
  $|V^-| = |V^+|$ (Lemma~\ref{lem:card-bounds}) gives
  $|\mathrm{Verts}(H)| = 2|V^+|$.
\end{proof}

\begin{theorem}[Generic clique induction]
  \label{thm:clique-induction}
  \uses{thm:spanner-of-two-hop-dismountable,def:spanner}
  \lean{TemporalGraph.Subgraph.exists_isSpannerOf_ncard_le_of_forall_not_twoHopDismountable}
  \leanok
  Let $G$ be symmetric and loopless, and $f : \mathbb{N} \to \mathbb{N}$ with
  $f(n) + 4 \le f(n + 1)$ for all $n$. If every finite, nontrivial,
  simple, proper, symmetric clique that is not $\{1,2\}$-hop dismountable has a
  symmetric spanner with at most $f(n)$ edges, $n$ its number of vertices,
  then so does every finite, simple, proper, symmetric clique.
\end{theorem}

\begin{proof}
  \leanok
  \uses{lem:dismount-clique,lem:mount-size}
  By strong induction on $n$; with at most one vertex, $H$ has no edge. If $H$
  is $\{1,2\}$-hop dismountable, dismount a vertex, recurse, and mount it back
  with at most $4$ edges (Theorem~\ref{thm:spanner-of-two-hop-dismountable},
  the paper's Observation~19, that is SAND 2025, Theorem~2.1 and \S~2.1):
  $f(n) + 4 \le f(n + 1)$. Otherwise the hypothesis applies. This is the
  induction that proves both the $O(n \log n)$ bound
  (Theorem~\ref{thm:clique-spanner-ncard-le}) and the linear bound below.
\end{proof}

\begin{theorem}[Temporal cliques admit spanners of size $7n$ --- Baligács Theorem 2]
  \label{thm:clique-spanner-seven-mul}
  \uses{thm:clique-induction,thm:spanner-of-not-two-hop-linear}
  \lean{TemporalGraph.Subgraph.exists_isSpannerOf_ncard_le_seven_mul}
  \leanok
  Let $G$ be symmetric and loopless. Every finite, simple, proper, symmetric
  clique $H$ on $n$ vertices has a symmetric spanner with at most $7n$ edges.
\end{theorem}

\begin{proof}
  \leanok
  Theorem~\ref{thm:clique-induction} with $f(n) = 7n$, which satisfies
  $7n + 4 \le 7(n + 1)$, and Theorem~\ref{thm:spanner-of-not-two-hop-linear} for
  the non-dismountable case.
\end{proof}

\begin{theorem}[Linear spanners of temporal cliques --- Baligács Theorem 2, answering Question 1]
  \label{thm:linear-main}
  \uses{thm:clique-spanner-seven-mul}
  \lean{TemporalGraph.Subgraph.exists_isBigO_isSpannerOf_linear}
  \leanok
  There is a function $f(n) = O(n)$, chosen once for all temporal cliques, such
  that every finite, simple, proper temporal clique $H$ on $n$ vertices (with
  $H$ and its ambient graph symmetric, the ambient graph loopless) admits a
  symmetric spanner with at most $f(n)$ edges. One can take $f(n) = 7n$. This
  improves Theorem~\ref{thm:main-spanner-bound}.
\end{theorem}

\begin{proof}
  \leanok
  Take $f(n) = 7n$, which bounds the spanners of every such clique by
  Theorem~\ref{thm:clique-spanner-seven-mul}, and is trivially $O(n)$.
\end{proof}
